% Source: https://artofproblemsolving.com/community/c3h2256948_ellipses_and_rhombi
% Stable identity: aops-topic-2256948-q01
Let the ellipse have centre $O$, semi-major axis $a$ and semi-minor axis $b$; then $a=\dfrac{24}{2}=12$. The axes are perpendicular, so each side of the rhombus is the hypotenuse of a right triangle with legs $a$ and $b$.

\begin{asy}
usepackage("amsmath");
size(7cm);
defaultpen(fontsize(10pt));
pair O=(0,0);
draw(ellipse(O,12,5),blue+linewidth(0.9));
draw((12,0)--(0,5)--(-12,0)--(0,-5)--cycle,linewidth(0.9));
draw((-12,0)--(12,0),dashed); draw((0,-5)--(0,5),dashed);
draw((0.6,0)--(0.6,0.6)--(0,0.6));
dot(O); label("$O$",O,dir(-135));
label("$12$",(6,0),S); label("$b$",(0,2.5),W); label("$13$",(6,2.5),NE);
\end{asy}

Hence $a^2+b^2=13^2$, so $b^2=169-144=25$ and $b=5$. The area of an ellipse with semi-axes $a$ and $b$ is $\pi ab$, so the area is
\[\pi\cdot12\cdot5=\boxed{60\pi}.\]
